% =====================================================================
%  Seance 6 -- C. Cours 3 : gouvernance + synthese de la mission (18--22)
% =====================================================================
\section{Gouvernance et clôture de la mission NOSYS}

\begin{frame}{Gouvernance 1 --- la responsabilité d'une alerte}
\small
\begin{columns}[t]
\begin{column}{0.5\textwidth}
\footnotesize\begin{block}{Deux scénarios qui arrivent tous les mois}
\begin{itemize}
\item \textbf{Alerte fausse traitée comme vraie} : 3 h perdues. Coût : fatigue, perte de confiance --- la « fatigue d'alertes ».
\item \textbf{Alerte vraie négligée} : l'incident du 22/09 se produit sans
intervention. Coût : indisponibilité, SLA violé.
\end{itemize}
\end{block}
\begin{wbox}[Le principe à retenir]
Le modèle \textbf{propose}, l'humain \textbf{dispose}. La responsabilité
d'une action (couper un service, appeler l'astreinte) reste
\textbf{humaine} --- et le seuil choisi (séance 3) est une décision
d'organisation, pas un réglage technique neutre.
\end{wbox}
\end{column}
\begin{column}{0.46\textwidth}
\begin{ebox}[Ce qu'une bonne organisation documente]
\begin{itemize}
\item Qui a fixé le seuil, quand, sur quelles métriques de validation.
\item Que faire d'une alerte : procédure écrite (runbook), délai cible.
\item Le suivi des faux positifs : chaque alerte fausse est
\textbf{étiquetée} et nourrit le prochain retraining.
\end{itemize}
\end{ebox}
\end{column}
\end{columns}
\end{frame}

\begin{frame}{Gouvernance 2 --- la model card : le modèle a une notice}
\begin{center}\footnotesize
\begin{tabular}{@{}p{3.2cm}p{9.2cm}@{}}
\toprule
\textbf{Rubrique} & \textbf{Contenu pour le modèle NOSYS}\\
\midrule
Objet & prédire une incident sur les 30 min suivantes, par fenêtre de 10 min
et par serveur\\
\addlinespace[1pt]
Données & \file{nosys\_features.csv}, 34 560 fenêtres, 8 serveurs, 30 jours
(split temporel au 22/09)\\
\addlinespace[1pt]
Performances & AUC test \textbf{0,80} ; AUC train 1,00 --- signe de marge de
manoeuvre à surveiller\\
\addlinespace[1pt]
Limites & incidents rares (0,35 \%) ; calibré sur un seul mois ; les
fenêtres sans incident ne permettent pas de mesurer l'AUC\\
\addlinespace[1pt]
Usage prévu & appuyer la décision du NOC --- jamais remplacer la décision
humaine\\
\bottomrule
\end{tabular}
\end{center}
\vspace{1mm}
\begin{kbox}[Une page, versionnée, à côté du modèle]
La model card (Mitchell et al., 2019) est le complément obligatoire du
fichier \file{.joblib} dans le registre. Si vous ne pouvez pas remplir la
rubrique « limites », vous ne connaissez pas assez votre modèle pour le
déployer.
\end{kbox}
\end{frame}

\begin{frame}{Gouvernance 3 --- vie privée des logs et auditabilité}
\footnotesize%
\begin{columns}[t]
\begin{column}{0.5\textwidth}
\begin{block}{Les logs parlent des personnes}
Les logs d'authentification contiennent des identifiants, des adresses IP,
des postes. Avant tout partage (LLM externe, prestataire, archive longue) :
\begin{itemize}
\item \textbf{pseudonymisation} : remplacer l'identifiant par une référence
réversible gardée à part ;
\item \textbf{minimisation} : ne garder que les champs utiles à l'analyse
(l'heure et le résultat d'une tentative, pas le contenu de la session) ;
\item durée de conservation définie et appliquée.
\end{itemize}
\end{block}
\end{column}
\begin{column}{0.46\textwidth}
\begin{block}{Auditabilité}
Pouvoir \textbf{remonter} : quelle version du modèle, sur quelles données,
a produit telle alerte à telle heure. Nos trois clés :
\begin{itemize}
\item registre de modèles (version + hash),
\item graine fixée (\textbf{5534}) $\Rightarrow$ données régénérables,
\item journal des prédictions conservé.
\end{itemize}
\end{block}
\end{column}
\end{columns}
\end{frame}

\begin{frame}{Synthèse de la mission NOSYS --- les 6 livrables de l'UE}
\begin{center}\footnotesize
\begin{tabular}{@{}p{0.7cm}p{4.0cm}p{4.4cm}p{3.2cm}@{}}
\toprule
\textbf{\#} & \textbf{Livrable} & \textbf{Résultat clé} & \textbf{Statut}\\
\midrule
1 & profil et saisonnalité des serveurs & cycle jour/nuit, incident du 09/09 & prêts\\
\addlinespace[1pt]
2 & \file{nosys\_features.csv} supervisé & 34 560 fenêtres, 120 positives
(0,35 \%), 0 fuite & prêts\\
\addlinespace[1pt]
3 & anomalies + classification & forêt : AUC test \textbf{0,80}, seuil calibré & prêts\\
\addlinespace[1pt]
4 & SARIMA + prédiction de panne & rappel 99,94 \% si 70 \% des incidents
évités & prêts\\
\addlinespace[1pt]
5 & analyse des logs, intrusion du 5/10 & brute force \file{bdd-principal}, 125 échecs & prêts\\
\addlinespace[1pt]
6 & industrialisation (cette séance) & pipeline joblib, latence, fiche MLOps & en TP\\
\bottomrule
\end{tabular}
\end{center}
\vspace{1mm}
\small\begin{kbox}[Ce qui est prêt pour la production]
La chaîne \textbf{features $\to$ forêt $\to$ scoring} est
reproductible de bout en bout (graine fixée, split temporel, modèle
sauvegardé). Ce qui manque avant production réelle : un mois supplémentaire
de recul, la supervision du service, et la procédure de rollback du B3.
\end{kbox}
\end{frame}
